% Chapter 14 draft for textbook inclusion. Assumes a textbook preamble with amsmath, amssymb, array, booktabs, graphicx, hyperref, and verbatim.
\chapter[Likelihood, GLMs, and Hypothesis Testing]{Likelihood, GLMs, and Hypothesis Testing Beyond Least Squares}
\label{chap:likelihood-glms-hypothesis-testing}

\section{The Big Picture}

Least squares gave us a way to fit a regression model by minimizing squared residuals.  Chapter~13 used the Gaussian-error theory behind least squares to build coefficient confidence intervals, mean-response intervals, and prediction intervals.

This chapter asks a different question:
\begin{quote}
  What if we start with a probability model for the response and fit the unknown parameters by choosing the values that make the observed data most plausible?
\end{quote}
That idea is called \emph{maximum likelihood}.  It explains why ordinary least squares is the natural estimator under Gaussian errors, and it gives us a way to build models when least squares is not appropriate, such as logistic regression for binary responses and Poisson regression for counts.

The goal is not to turn this chapter into a complete generalized-linear-model course.  The goal is to give the statistical logic needed to recognize likelihood-based models, understand their assumptions, and interpret their tests carefully.

\paragraph{Connection to earlier chapters.}
Chapters~5--10 built the linear-regression model, its assumptions, and the distributional behavior of \(\widehat{\boldsymbol{\beta}}\).  Chapter~12 emphasized diagnostics.  Chapter~13 turned Gaussian fixed-design theory into interval estimates.  This chapter explains the likelihood logic that sits underneath those results and extends the same modeling mindset beyond Gaussian least squares.

\paragraph{Math Foundations Connection.}
Likelihood models use functions, optimization, gradients, and sometimes second-derivative checks.  Regression likelihoods also rely on matrix-vector multiplication, least-squares objectives, full-rank inverse formulas, and quadratic residual forms.  These ideas are supported by the Math Foundations textbook, especially Chapters~5, 7, 10, 12, and~13.

\section{Likelihood and Log-Likelihood}

Suppose we observe data \(Y_1,\ldots,Y_n\) from a probability model with parameter vector \(\boldsymbol{\theta}\).  For each possible value of \(\boldsymbol{\theta}\), the model assigns a probability density or probability mass to the observed data.

For independent observations, the likelihood is
\begin{equation}
  L(\boldsymbol{\theta}\mid y_1,\ldots,y_n)
  = \prod_{i=1}^n f(y_i\mid \boldsymbol{\theta}),
  \label{eq:ch14-likelihood}
\end{equation}
where \(f\) is the density or mass function of the response model.

The \emph{maximum likelihood estimator}, or MLE, is
\begin{equation}
  \widehat{\boldsymbol{\theta}}_{MLE}
  =\arg\max_{\boldsymbol{\theta}\in\Theta}
  L(\boldsymbol{\theta}\mid y_1,\ldots,y_n).
\end{equation}
In words, the MLE is the parameter value that makes the observed data most likely under the model.

Because products are often awkward, we usually maximize the log-likelihood
\begin{equation}
  \ell(\boldsymbol{\theta})
  =\log L(\boldsymbol{\theta}\mid y_1,\ldots,y_n)
  =\sum_{i=1}^n \log f(y_i\mid \boldsymbol{\theta}).
  \label{eq:ch14-loglikelihood}
\end{equation}
The logarithm is increasing, so the value of \(\boldsymbol{\theta}\) that maximizes the likelihood also maximizes the log-likelihood.

\paragraph{Important interpretation.}
The likelihood is a function of the parameter after the data have been observed.  It is not the probability that the model is true.  It is a way to compare parameter values within an assumed model family.

\section{A First MLE Example: The Poisson Mean}

Let \(Y_1,\ldots,Y_n\) be independent observations from a Poisson distribution with mean \(\mu>0\):
\begin{equation}
  Y_i\sim \operatorname{Pois}(\mu).
\end{equation}
The probability mass function is
\begin{equation}
  f(y_i\mid \mu)=\frac{e^{-\mu}\mu^{y_i}}{y_i!}.
\end{equation}
The log-likelihood is
\begin{align}
  \ell(\mu)
  &=\sum_{i=1}^n \log\left(\frac{e^{-\mu}\mu^{y_i}}{y_i!}\right) \\
  &=\sum_{i=1}^n \left[-\mu+y_i\log(\mu)-\log(y_i!)\right].
\end{align}
Taking the derivative and setting it equal to zero gives
\begin{equation}
  0=-n+\frac{1}{\widehat{\mu}_{MLE}}\sum_{i=1}^n y_i,
\end{equation}
so
\begin{equation}
  \boxed{\widehat{\mu}_{MLE}=\bar{y}.}
\end{equation}
When at least one observed count is positive, this is the interior maximizer.  The all-zero sample is a boundary case: the likelihood is maximized at \(\mu=0\) if that boundary is allowed.
For a Poisson mean, maximum likelihood gives the sample mean.  This example is simple, but it illustrates the pattern: write a probability model, form the likelihood, and choose the parameter value that maximizes it.

\section{Gaussian-Error Regression and OLS as MLE}

Now return to the fixed-design linear model.  Under the strongest regression assumption,
\begin{equation}
  \boldsymbol{\epsilon}\mid\mathbf{X}
  \sim N(\mathbf{0},\sigma^2\mathbf{I}_n),
\end{equation}
so
\begin{equation}
  \mathbf{Y}\mid\mathbf{X}
  \sim N(\mathbf{X}\boldsymbol{\beta},\sigma^2\mathbf{I}_n).
  \label{eq:ch14-gaussian-linear-model}
\end{equation}
The unknown parameters are \(\boldsymbol{\beta}\) and \(\sigma^2\).

For observed response vector \(\mathbf{y}\), the log-likelihood has the form
\begin{equation}
  \ell(\boldsymbol{\beta},\sigma^2)
  = -\frac{n}{2}\log(\sigma^2)
    -\frac{1}{2\sigma^2}(\mathbf{y}-\mathbf{X}\boldsymbol{\beta})^T
    (\mathbf{y}-\mathbf{X}\boldsymbol{\beta})
    + C,
  \label{eq:ch14-gaussian-loglik}
\end{equation}
where \(C\) does not depend on \(\boldsymbol{\beta}\) or \(\sigma^2\).

For fixed \(\sigma^2\), maximizing \eqref{eq:ch14-gaussian-loglik} over \(\boldsymbol{\beta}\) is equivalent to minimizing
\begin{equation}
  RSS(\mathbf{b})=(\mathbf{y}-\mathbf{X}\mathbf{b})^T(\mathbf{y}-\mathbf{X}\mathbf{b}).
\end{equation}
Therefore, under Gaussian errors, the maximum likelihood estimator for \(\boldsymbol{\beta}\) is the ordinary least-squares estimator.  If \(\mathbf{X}\) has full column rank, then
\begin{equation}
  \boxed{
  \widehat{\boldsymbol{\beta}}_{MLE}
  =\widehat{\boldsymbol{\beta}}
  = (\mathbf{X}^T\mathbf{X})^{-1}\mathbf{X}^T\mathbf{y}.
  }
\end{equation}

This is the key bridge: least squares is an algebraic projection method, but under Gaussian errors it is also a likelihood method.

\section[Estimating sigma squared: MLE Versus Residual Variance]{Estimating \(\sigma^2\): MLE Versus Residual Variance}

The MLE for \(\sigma^2\) is
\begin{equation}
  \widehat{\sigma}^2_{MLE}
  =\frac{1}{n}(\mathbf{y}-\mathbf{X}\widehat{\boldsymbol{\beta}})^T
  (\mathbf{y}-\mathbf{X}\widehat{\boldsymbol{\beta}})
  =\frac{RSS}{n}.
  \label{eq:ch14-sigma2-mle}
\end{equation}
This is the empirical mean squared residual.

However, after fitting \(q=p+1\) coefficients, the residuals have only \(n-q\) degrees of freedom.  The usual unbiased residual variance estimator is
\begin{equation}
  s^2
  =\frac{RSS}{n-q}
  =\frac{RSS}{n-p-1}.
  \label{eq:ch14-sigma2-unbiased}
\end{equation}
This is the quantity used in the standard errors, \(t\)-tests, and intervals from the previous chapters.

The distinction matters:
\begin{center}
\begin{tabular}{p{0.28\textwidth}p{0.30\textwidth}p{0.32\textwidth}}
\toprule
\textbf{Quantity} & \textbf{Formula} & \textbf{Main role} \\
\midrule
MLE for \(\sigma^2\) & \(RSS/n\) & likelihood maximization \\
Unbiased residual variance & \(RSS/(n-q)\) & standard errors and classical inference \\
\bottomrule
\end{tabular}
\end{center}

The same issue appears in the sample-variance formula: using \(1/n\) gives the MLE under a normal model, while using \(1/(n-1)\) gives the usual unbiased sample variance.

\section{Why the Assumptions Matter}

Maximum likelihood always starts with a probability model.  If the model changes, the likelihood changes, and the interpretation of the estimator changes.

For ordinary least squares, the formula
\begin{equation}
  (\mathbf{X}^T\mathbf{X})^{-1}\mathbf{X}^T\mathbf{y}
\end{equation}

enjoys several roles:
\begin{itemize}
  \item it minimizes squared residuals;
  \item it is the orthogonal-projection estimator from the design matrix viewpoint;
  \item it is the MLE for \(\boldsymbol{\beta}\) under Gaussian independent errors.
\end{itemize}
If the Gaussian-error assumption fails, the estimator may still be the least-squares estimator, but the exact Gaussian likelihood and exact small-sample \(t\)- and \(F\)-based conclusions no longer follow automatically.

This is why diagnostics still matter.  A software package can print standard errors, confidence intervals, and p-values even when the model assumptions are questionable.  The analyst is responsible for deciding whether those summaries are defensible.

\section{Generalized Linear Models: The Three-Part Idea}

Linear regression assumes a quantitative response whose conditional mean is linear in the predictors and whose errors are Gaussian.  Many data-analysis problems do not fit that pattern.

A \emph{generalized linear model}, or GLM, keeps the idea of a linear predictor but changes the response distribution and the way the mean connects to the predictors.  The basic pieces are:
\begin{enumerate}
  \item a response distribution, such as Bernoulli for binary data or Poisson for counts;
  \item a linear predictor \(\eta=\mathbf{x}^T\boldsymbol{\beta}\);
  \item a link function connecting the conditional mean \(\mu(\mathbf{x})=\mathbb{E}[Y\mid \mathbf{x}]\) to \(\eta\).
\end{enumerate}
In symbols, the link function satisfies
\begin{equation}
  g(\mu(\mathbf{x}))=\mathbf{x}^T\boldsymbol{\beta}.
\end{equation}

The details of GLM theory are beyond this course.  The important idea is that likelihood lets us fit models where the response is not naturally Gaussian and where a straight-line mean model would violate basic constraints.

\section{Logistic Regression for Binary Responses}

Suppose \(Y\) is binary, with \(Y=1\) for success and \(Y=0\) for failure.  A linear model for the probability,
\begin{equation}
  P(Y=1\mid \mathbf{x})=\mathbf{x}^T\boldsymbol{\beta},
\end{equation}
would be dangerous because \(\mathbf{x}^T\boldsymbol{\beta}\) can be negative or greater than one.

Logistic regression avoids this problem by using the logistic function:
\begin{equation}
  P(Y=1\mid \mathbf{x})
  =\frac{\exp(\mathbf{x}^T\boldsymbol{\beta})}
  {1+\exp(\mathbf{x}^T\boldsymbol{\beta})}.
  \label{eq:ch14-logistic-prob}
\end{equation}
This quantity is always between zero and one.

Equivalently,
\begin{equation}
  \log\left(\frac{P(Y=1\mid \mathbf{x})}{1-P(Y=1\mid \mathbf{x})}\right)
  =\mathbf{x}^T\boldsymbol{\beta}.
\end{equation}
The left side is the log-odds.  Thus logistic-regression coefficients are linear effects on log-odds, not direct additive effects on probability.

\paragraph{Classification rule.}
Logistic regression estimates a probability.  To turn that probability into a class label, the analyst chooses a threshold, such as classifying as \(Y=1\) when \(\widehat{P}(Y=1\mid\mathbf{x})>0.5\).  That threshold is a decision rule, not a law of the model.

\paragraph{Scope warning.}
This section introduces logistic regression as a likelihood-based model.  Broader classification metrics and other classification algorithms are separate topics and are not developed here.

\section{Poisson Regression for Count Responses}

Suppose \(Y\) is a count: number of arrivals, number of defects, number of events, or number of items observed in a fixed window.  A Poisson model is often a first candidate:
\begin{equation}
  Y\mid\mathbf{x}\sim \operatorname{Pois}(\mu(\mathbf{x})).
\end{equation}
The mean \(\mu(\mathbf{x})\) must be positive.  A direct linear mean \(\mu(\mathbf{x})=\mathbf{x}^T\boldsymbol{\beta}\) can be negative, so it is usually not a safe final model.

Poisson regression uses a log link:
\begin{equation}
  \log \mu(\mathbf{x})=\mathbf{x}^T\boldsymbol{\beta},
  \qquad\text{or equivalently}\qquad
  \mu(\mathbf{x})=\exp(\mathbf{x}^T\boldsymbol{\beta}).
  \label{eq:ch14-poisson-log-link}
\end{equation}
This guarantees a positive fitted mean.

A coefficient in a Poisson log-link model has a multiplicative interpretation.  Increasing predictor \(x_j\) by one unit changes the fitted mean by a factor of \(\exp(\beta_j)\), holding the other predictors fixed.

\section{Hypothesis Testing Revisited}

Earlier chapters introduced two common tests in Gaussian linear regression:
\begin{center}
\begin{tabular}{p{0.22\textwidth}p{0.34\textwidth}p{0.30\textwidth}}
\toprule
\textbf{Test} & \textbf{Typical question} & \textbf{Main setting} \\
\midrule
\(t\)-test & Is one coefficient different from zero? & one coefficient at a time \\
\(F\)-test & Does a larger linear model improve over a smaller nested linear model? & Gaussian linear models \\
\bottomrule
\end{tabular}
\end{center}

In this framework, ANOVA can be viewed as a special case of comparing nested Gaussian linear models, often when the predictors are categorical group indicators.

This chapter does not re-derive those tests.  The key point is that they rely on the Gaussian regression setup, the residual variance estimate, and the distributional results developed earlier.

Likelihood gives a more general model-comparison idea: compare how well a smaller model and a larger model explain the observed data through their maximized log-likelihoods.

\section{Likelihood-Ratio Tests}

Suppose we have two nested likelihood-based models:
\begin{itemize}
  \item a smaller model, with fewer free parameters;
  \item a larger model, which contains the smaller model as a special case.
\end{itemize}
Let \(\ell_{small}\) and \(\ell_{full}\) be the maximized log-likelihoods for the two models.  Since the full model has at least as much flexibility,
\begin{equation}
  \ell_{full}\ge \ell_{small}.
\end{equation}
The likelihood-ratio test statistic is
\begin{equation}
  \boxed{
  \Lambda = 2\left(\ell_{full}-\ell_{small}\right).
  }
  \label{eq:ch14-lrt}
\end{equation}
Under regularity conditions and under the null hypothesis that the smaller model is sufficient, the large-sample reference distribution is
\begin{equation}
  \Lambda \approx \chi^2_{d},
\end{equation}
where \(d\) is the difference in the number of free parameters between the full and smaller models.

\paragraph{Interpretation.}
A large likelihood-ratio statistic means that the full model improved the maximized log-likelihood substantially.  That is evidence against the smaller model, provided the models are nested and the likelihood assumptions are appropriate.

\paragraph{Connection to the \(F\)-test.}
For Gaussian linear models, likelihood-ratio tests and \(F\)-tests are related because both compare nested models using residual error information.  They are not the same object, and the likelihood-ratio test is more general: it applies to many maximum-likelihood models, including GLMs.

\section{Common Mistakes}

\subsection*{Mistake 1: Treating likelihood as the probability the model is true}

Likelihood compares parameter values inside a model family after observing the data.  It is not the probability that the model itself is true.

\subsection*{Mistake 2: Using a linear mean where the response has constraints}

Probabilities must lie between zero and one.  Poisson means must be positive.  Link functions help enforce these constraints.

\subsection*{Mistake 3: Forgetting the assumptions behind classical tests}

The usual \(t\)- and \(F\)-tests depend on the Gaussian-error regression framework.  Diagnostics and modeling judgment still matter.

\subsection*{Mistake 4: Comparing non-nested models with a likelihood-ratio test}

The likelihood-ratio test described here is for nested models.  If two models are not nested, use a different comparison tool, such as validation error or an information criterion.

\subsection*{Mistake 5: Turning logistic regression into a black-box classifier too quickly}

Logistic regression estimates probabilities.  Classification requires a threshold and a decision context.  The threshold should be chosen to match the operational problem.

\section{Chapter Summary}

Maximum likelihood fits a probabilistic model by choosing the parameter values that maximize the likelihood of the observed data.  For independent data,
\begin{equation}
  L(\boldsymbol{\theta})=\prod_{i=1}^n f(y_i\mid\boldsymbol{\theta}),
  \qquad
  \ell(\boldsymbol{\theta})=\sum_{i=1}^n\log f(y_i\mid\boldsymbol{\theta}).
\end{equation}

Under Gaussian errors,
\begin{equation}
  \mathbf{Y}\mid\mathbf{X}\sim N(\mathbf{X}\boldsymbol{\beta},\sigma^2\mathbf{I}_n),
\end{equation}
so the MLE for \(\boldsymbol{\beta}\) is the ordinary least-squares estimator.  The MLE for \(\sigma^2\) is \(RSS/n\), while the usual unbiased residual variance estimator is \(RSS/(n-p-1)\).

GLMs extend the regression idea by combining a response distribution, a linear predictor, and a link function.  Logistic regression models binary responses with probabilities between zero and one.  Poisson regression models count responses with positive means.

Likelihood-ratio tests compare nested likelihood-based models using the improvement in maximized log-likelihood.  They are powerful and general, but they require nested models, appropriate assumptions, and careful interpretation.

\section{Practice and Workflow}

\subsection*{Focused Study Checklist}

Use this checklist after studying likelihood and GLMs.

\begin{enumerate}
  \item Can you explain the difference between a likelihood and a probability statement about a model?
  \item Can you write the likelihood for independent observations as a product?
  \item Can you explain why maximizing log-likelihood gives the same estimator as maximizing likelihood?
  \item Can you show why Gaussian-error least squares is also maximum likelihood for \(\boldsymbol{\beta}\)?
  \item Can you distinguish \(RSS/n\) from \(RSS/(n-p-1)\)?
  \item Can you explain why logistic regression keeps fitted probabilities between zero and one?
  \item Can you explain why Poisson regression usually uses a log link?
  \item Can you state when a likelihood-ratio test is appropriate?
  \item Can you explain why diagnostics still matter before trusting likelihood-based inference?
\end{enumerate}

\paragraph{Coding companion reference.}
Package-specific commands for fitting logistic regression, Poisson regression, GLMs, and likelihood-ratio tests belong in the separate Coding and Software Cookbook companion.  This chapter keeps the statistical meaning and model structure.
